\documentclass[10pt,a4paper]{amsart}
\usepackage[margin=19mm]{geometry}
\usepackage{amsmath,amssymb,mathtools}
\usepackage[scr=rsfs,cal=euler]{mathalfa}
\usepackage{xfrac}
\usepackage[unicode,hidelinks,bookmarks=false]{hyperref}
\newcommand{\ie}{i.e.\ }
\newcommand{\cf}{cf.\ }
\newcommand{\resp}{respectively\ }
\newcommand{\Subsection}{\subsection}
\providecommand{\nobreakdash}{\nobreak\hbox{-}}
\setlength{\emergencystretch}{2em}
\multlinegap0pt
\usepackage{mathtools}
\usepackage{amssymb}
\usepackage[shortlabels]{enumitem}
\usepackage{tikz}
\newtheorem{proposition}{Proposition}
\newtheorem{corollary}{Corollary}
\usepackage{cleveref}
\crefname{proposition}{Proposition}{Propositions}
\crefname{corollary}{Corollary}{Corollarys}
\newcommand{\R}{\mathbb{R}}
\newcommand{\dd}{\mathrm{dist}}
\newcommand{\dist}{{\mathrm{dist}_Y|_X}}
\title{Morse Theory of Euclidean Distance Functions from Algebraic Hypersurfaces}
\author{Andrea Guidolin, Antonio Lerario, Isaac Ren, Martina Scolamiero}
\date{Source: arXiv:2402.08639v4 (CC BY 4.0)}
\begin{document}
\maketitle

\noindent\emph{Source and licence.} Morse Theory of Euclidean Distance Functions from Algebraic Hypersurfaces. Version arXiv:2402.08639v4 (CC BY 4.0), distributed by arXiv under the Creative Commons Attribution 4.0 International licence, http://creativecommons.org/licenses/by/4.0/, as recorded in the arXiv licence metadata for this exact version and shown on the abstract page https://arxiv.org/abs/2402.08639v4. Full licence notice: \texttt{licenses/arxiv-2402.08639-CC-BY-4.0.txt}.\par\smallskip

\noindent\textit{Editorial scope.} Counted unit: the article's proposition propo:subgradient2 (source0.tex lines 589--604). The article's complete original proof of it is delivered with it (source0.tex lines 606--695, one \texttt{proof} environment of 90 lines). No mathematics is written, repaired or re-derived: the statement and the proof are the article's own lines, in the article's own notation, and no retained line is substituted or re-flowed.  Retained uncounted as the source-local context the proof needs: lines 452--469; lines 477--486; lines 545--555.  Nothing was omitted from any retained span, so the package carries no omission record.  No retained line contains a citation, so the wrapper carries no bibliography and no bibliography entry.  No retained line contains a figure environment, an image-inclusion command or drawing code, so no image is delivered and the package declares no asset dependency.  Editorial formatting only, in addition: an AMS \texttt{amsart} preamble that restates the article's own declarations and macros used by the retained lines, namely the environments \texttt{proposition}, \texttt{corollary} on separate counters, together with the standard packages those lines need.  The retained environments print in this document as Proposition 1, Corollary 1 in order of appearance, whereas the article prints the same items with the numbering of its own full document.  The complete native source of the article as distributed in the arXiv e-print for this exact version is bundled unchanged as \texttt{sources/154502/source0.tex}.
\begin{proposition}[Danskin's theorem]\label{P:danskin}
Let $W \subseteq \R^m$ be compact, and let $\phi \colon \R^n \times W \to \R$ be continuous. Assume that for every $w \in W$ the function $\phi(-,w)\colon \R^n \to \R$ is differentiable, and that for every $x \in \R^n$ the map
\[
w \mapsto \nabla_x \phi(x,w)
\]
is continuous on $W$. Define
\[
f(x) \coloneqq \max_{w \in W} \phi(x,w).
\]
Then
\[
\partial_x f = \mathrm{co}\left\{ \nabla_x \phi(x,z) \;\middle|\; z \in Z(x)\right\},
\]
where
\[
W(x) \coloneqq \left\{ w \in W \;\middle|\; \phi(x,w)=\max_{w \in W}\phi(x,w)\right\}.
\]
\end{proposition}

\begin{proposition}\label{P:squared-distance-subgradient}
Let $Y \subseteq \R^n$ be closed, and define
\[
\eta_Y(x)\coloneqq \frac{1}{2}\dd_Y(x)^2.
\]
Then, for every $x \in \R^n$,
\[
\partial_x(\eta_Y)= \mathrm{co}\left\{ x-y \;\middle|\; y \in B(x,\dd_Y(x)) \cap Y \right\}.
\]
\end{proposition}

\begin{corollary}\label{propo:subgradient1}
Let $Y\subseteq \R^n$ be closed and let $x\in \R^n \setminus Y$. Then
\[
\partial_x(\dd_Y) = \mathrm{co}\left\{
\left.
\frac{x-y}{\|x-y\|}
\,\right|\,
y \in B(x,\dd_Y(x)) \cap Y
\right\}.
\]
\end{corollary}

\begin{proposition}\label{propo:subgradient2}
Let $Y \subseteq \R^n$ be closed and let $X \subseteq \R^n$ be a smooth manifold. Then, for every $x \in X \setminus Y$,
\[
\partial_x(\dist)=\mathrm{proj}_{T_xX}\bigl(\partial_x(\dd_Y)\bigr),
\]
where $\mathrm{proj}_{T_xX}\colon \R^n \to T_xX$ denotes the orthogonal projection. More explicitly,
\[
\partial_x(\dist)
=
\mathrm{co}\left\{
\frac{\mathrm{proj}_{T_xX}(x-y)}{\|x-y\|}
\;\middle|\;
y \in B(x,\dd_Y(x)) \cap Y
\right\}.
\]
\end{proposition}

\begin{proof}
Fix $x_0 \in X \setminus Y$, and choose a smooth chart
\[
\psi \colon U \subseteq \R^{\dim X} \to X
\]
such that $\psi(0)=x_0$.

As in the proof of \cref{P:squared-distance-subgradient}, after shrinking $U$ if necessary there exists $R>0$ such that
\[
B(\psi(u),\dd_Y(\psi(u))) \cap Y \subseteq B(0,R)
\qquad \text{for every } u \in U.
\]
Set
\[
W \coloneqq Y \cap B(0,R)
\]
and define
\[
F(u)\coloneqq -\frac{1}{2}\dd_Y|_X(\psi(u))^2.
\]
Then
\[
F(u)=\max_{y \in W}\Phi(u,y),
\]
where
\[
\Phi(u,y)\coloneqq
\langle \psi(u),y\rangle-\frac{\|\psi(u)\|^2}{2}-\frac{\|y\|^2}{2}.
\]
The function $\Phi$ is continuous, for every $y \in Z$ the function $\Phi(-,y)$ is smooth, and
\[
\nabla_u\Phi(u,y)=D\psi(u)^{\top}(y-\psi(u))
\]
depends continuously on $y$. Therefore \cref{P:danskin} applies and gives
\[
\label{Eq:F-subdif-first}
\partial_0F
=
\mathrm{co}\left\{
D\psi(0)^{\top}(y-x_0)
\;\middle|\;
y \in B(x_0,\dd_Y(x_0)) \cap Y
\right\}.
\]
Moreover, by the chain rule, we have
\[
\label{Eq:F-subdif-second}
\partial_0 F = D \psi(0)^\top \partial_{x_0} \left( -\frac{1}{2} \dist^2 \right).
\]
Identifying $T_{x_0} X$ with the image of $D \psi(0)$, Equations \eqref{Eq:F-subdif-first} and \eqref{Eq:F-subdif-second} give the following equality between projections onto $T_{x_0} X$:
\[
\operatorname{proj}_{T_{x_0}} \partial_{x_0} \left( -\frac{1}{2} \dist^2 \right) = \mathrm{co} \left\{
\mathrm{proj}_{T_{x_0}X} (y - x_0)
\;\middle|\;
y \in B(x_0, \dd_Y(x_0)) \cap Y
\right\}.
\]
The subdifferential $\partial_{x_0} \left( -\frac{1}{2} \dist^2 \right)$ is already in $T_{x_0} X$ by construction, so we can omit the projection on the lefthand side. Taking the opposite sign, we get
\[
\partial_{x_0}\left(\frac{1}{2}\dist^2\right)
=
\mathrm{co}\left\{
\mathrm{proj}_{T_{x_0}X}(x_0-y)
\;\middle|\;
y \in B(x_0,\dd_Y(x_0)) \cap Y
\right\}.
\]
Since $x_0 \notin Y$, one has $\dist(x_0)=\dd_Y(x_0)>0$, and therefore
\[
\partial_{x_0}(\dist)
=
\frac{1}{\dist(x_0)}
\partial_{x_0}\left(\frac{1}{2}\dist^2\right).
\]
Using that $\|x_0-y\|=\dd_Y(x_0)$ for every $y \in B(x_0,\dd_Y(x_0)) \cap Y$, we conclude that
\[
\partial_{x_0}(\dist)
=
\mathrm{co}\left\{
\frac{\mathrm{proj}_{T_{x_0}X}(x_0-y)}{\|x_0-y\|}
\;\middle|\;
y \in B(x_0,\dd_Y(x_0)) \cap Y
\right\}.
\]
By \cref{propo:subgradient1}, this is exactly
\[
\partial_{x_0}(\dist)=\mathrm{proj}_{T_{x_0}X}\bigl(\partial_{x_0}(\dd_Y)\bigr),
\]
because orthogonal projection is linear.
\end{proof}

\end{document}
